\chapter{Memahami Gaya Koneksi: Anchor, Island, Wave}

\begin{insightbox}
\textit{``Kunci hubungan yang sukses bukanlah menemukan pasangan sempurna, tapi memahami \highlight{bagaimana} kalian berdua terhubung. Ketika kamu mengenali gaya koneksimu dan pasanganmu, konflik bukan lagi personal --- itu hanya perbedaan wiring.''}
\end{insightbox}

\section{Apa itu Attachment Styles?}

\subsection{Sejarah Singkat}

Teori koneksi (attachment theory) pertama kali dikembangkan oleh psikolog John Bowlby pada tahun 1960-an. Bowlby mempelajari bagaimana bayi bereaksi saat dipisahkan dari caregiver-nya dan menemukan bahwa cara caregiver merespons kebutuhan bayi akan membentuk ``template'' hubungan untuk seumur hidup.

Dr. Stan Tatkin, dalam bukunya \textit{Wired for Dating}, menyederhanakan tiga dekade penelitian attachment menjadi tipe yang mudah dipahami: \highlight{Anchor} (Jangkar), \highlight{Island} (Pulau), dan \highlight{Wave} (Gelombang).

\begin{praktikbox}[Definisi Secure Functioning]
\keyterm{Secure functioning} adalah hubungan di mana:
\begin{itemize}
    \item Kalian saling menjaga keselamatan dan keamanan
    \item Sensitif terhadap kebutuhan satu sama lain
    \item Cepat memperbaiki kesalahan yang terjadi
    \item Berkolaborasi sebagai tim
    \item Menjalankan \textit{true mutuality} --- apa yang baik untukku juga baik untukmu
\end{itemize}
\end{praktikbox}

\section{Anchor (Jangkar): Dua Lebih Baik dari Satu}

\subsection{Siapa itu Anchor?}

Anchor adalah tipe yang paling \textit{secure} dalam hubungan. Mereka tumbuh dengan caregiver yang responsif, sehingga internalize bahwa dunia adalah tempat yang aman dan orang lain bisa diandalkan.

\begin{center}
\begin{tikzpicture}[
    anchorbox/.style={draw=sagegreen, line width=2pt, rounded corners=15pt, 
                   fill=sagegreen!10, minimum width=10cm, inner sep=15pt}
]
    \node[anchorbox] (box) {
        \begin{minipage}{9cm}
            \centering
            {\Large\bfseries\color{sagegreen} KARAKTERISTIK ANCHOR}\\[0.5cm]
            \begin{itemize}[leftmargin=*, itemsep=0.2cm]
                \item \textbf{Secure} --- merasa aman sebagai individu dan dalam hubungan
                \item \textbf{Willing to commit} --- tidak takut komitmen
                \item \textbf{Fully share} --- bisa berbagi secara penuh dengan pasangan
                \item \textbf{Generally happy} --- cenderung bahagia secara umum
                \item \textbf{Adaptable} --- fleksibel menyesuaikan kebutuhan pasangan
                \item \textbf{Collaborative} --- percaya ``dua lebih baik dari satu''
                \item \textbf{Affectionate} --- mudah menunjukkan kasih sayang
                \item Tidak takut \textit{abandonment} maupun \textit{engulfment}
            \end{itemize}
        \end{minipage}
    };
\end{tikzpicture}
\end{center}

\subsection{Anchor sebagai Anak}

Anchor berasal dari lingkungan keluarga yang:
\begin{enumerate}
    \item Menempatkan \textbf{hubungan di atas segalanya}
    \item Setidaknya satu caregiver yang \textbf{fully present} dan tersedia
    \item Caregiver \textbf{affectionate} --- suka memeluk, mencium, mengayun
    \item Menghabiskan waktu \textbf{face-to-face, eye-to-eye, skin-to-skin}
    \item Saat anak marah/sedih, caregiver \textbf{menenangkan dengan efektif}
    \item Tidak membeli anak dengan hadiah saat anak emosional
    \item Melihat anak secara akurat dan mengenalnya dengan dalam
    \item Menghargai \textbf{dependency} dan \textbf{autonomy} secara seimbang
\end{enumerate}

\begin{tipbox}[Bedah Perilaku Anchor]
\textbf{Sehari-hari:}
\begin{itemize}
    \item Mudah meminta bantuan dan menerima bantuan
    \item Tidak merasa terancam jika pasangan butuh ruang
    \item Tidak merasa cekik jika pasangan ingin dekat
    \item Bisa sendiri atau bersama pasangan dengan sama nyamannya
\end{itemize}

\textbf{Dalam Konflik:}
\begin{itemize}
    \item Mencari solusi win-win
    \item Tidak defensive secara berlebihan
    \item Tetap terhubung meski sedang marah
    \item Bisa mengatakan ``aku memilihmu'' tanpa ragu
\end{itemize}
\end{tipbox}

\subsection{Apakah Anchor Sempurna?}

\textbf{Tidak!} Anchor juga bisa:
\begin{itemize}
    \item Membuat lelucon yang tidak lucu
    \item Terlambat ke janji
    \item Lupa mematikan lampu
    \item Menyebalkan dalam hal-hal kecil
\end{itemize}

Tapi mereka \textbf{resilient} --- punya sumber daya internal yang kuat karena merasa \textit{tethered} (terikat) pada seseorang. Mereka tidak hidup sebagai \textit{lone wolves}.

\section{Island (Pulau): Aku Bisa Melakukannya Sendiri}

\subsection{Paradoks Island}

Island adalah tipe yang paling \textbf{kontradiktif}. Di permukaan, mereka tampak:
\begin{itemize}
    \item Mandiri
    \item Tidak butuh siapa-siapa
    \item Low maintenance
\end{itemize}

Tapi \textbf{paradoksnya}: Island seringkali adalah tipe yang \textbf{paling membutuhkan} pasangan, hanya saja mereka \textit{tidak tahu} cara menerima kebutuhan itu.

\begin{center}
\begin{tikzpicture}[
    island/.style={draw=primary, line width=2pt, rounded corners=15pt, 
                   fill=lightbg, minimum width=10cm, inner sep=15pt}
]
    \node[island] (box) {
        \begin{minipage}{9cm}
            \centering
            {\Large\bfseries\color{primary} KARAKTERISTIK ISLAND}\\[0.5cm]
            \begin{itemize}[leftmargin=*, itemsep=0.2cm]
                \item \textbf{Independent} --- sangat mandiri dan \textit{self-reliant}
                \item \textbf{Detail-oriented} --- perhatian terhadap detail tinggi
                \item \textbf{Logical \& Rational} --- berpikir logis dan rasional
                \item \textbf{Performance-focused} --- performa sangat penting
                \item \textbf{Appearance-conscious} --- peduli penampilan
                \item \textbf{Low maintenance} --- sepertinya tidak butuh banyak
                \item \textbf{Avoidant} --- menghindari saat stres (distancing group)
                \item \textbf{Self-soothing} --- menenangkan diri sendiri
            \end{itemize}
        \end{minipage}
    };
\end{tikzpicture}
\end{center}

\subsection{Island sebagai Anak}

Island tumbuh dalam keluarga yang:
\begin{enumerate}
    \item Menekankan \textbf{performa, kecerdasan, bakat, penampilan}
    \item \textbf{Menghambat dependency} --- anggap kebutuhan itu aib
    \item Caregiver bisa \textbf{dingin, meremehkan, atau dismissive}
    \item Saat anak sedih, caregiver memberi \textbf{uang/mainan} alih-alih perhatian
    \item Caregiver \textbf{menarik diri, aloof, atau tidak tersedia}
    \item Caregiver \textbf{mementingkan kebutuhan sendiri} di atas hubungan
    \item Saat anak ketakutan di malam hari, caregiver \textbf{tidak responsif}
    \item Saat anak kesakitan, caregiver \textbf{mengabaikan, mengkritik, atau meremehkan}
\end{enumerate}

\begin{warningbox}[Kenapa Island Tampak ``Baik-baik Saja''?]
Island anak seringkali tampak \textit{well-adjusted} karena:
\begin{itemize}
    \item Mereka bermain sendiri dengan bahagia
    \item Tidak menunjukkan kebutuhan secara terbuka
    \item Tampak ``dewasa'' untuk usianya
\end{itemize}

Tapi di dalam, mereka memiliki \textbf{self yang rapuh} yang tidak tahu cara menangani pikiran dan perasaan stres.
\end{warningbox}

\subsection{Mengapa Island Menjauh (Distancing)?}

Saat stres, Island masuk ke \keyterm{distancing group} --- menjauh dari orang lain. Ini karena:

\begin{enumerate}
    \item \textbf{Hubungan dekat = Stres:} Hubungan intim membawa terlalu banyak ancaman yang kabur
    \item \textbf{Self-soothing otomatis:} Refleks mereka adalah menenangkan diri sendiri, bukan mencari pasangan
    \item \textbf{Takut kehilangan otonomi:} Island sangat sensitif terhadap perasaan terjebak
    \item \textbf{Pengalaman keluarga:} Dari kecil, isolasi = relief, bukan kehilangan
\end{enumerate}

\begin{tipbox}[Deteksi Island dalam Hubungan]
\textbf{Warning signs:}
\begin{itemize}
    \item Sering lupa janji atau ``lupa'' memberi tahu kemana perginya
    \item Saat ada masalah, menghilang atau merenung sendiri
    \item Sulit bertransisi dari aktivitas solo ke interaksi dengan pasangan
    \item Saat dipeluk terlalu lama, mulai gelisah
    \item Saat ditanya ``Kamu kenapa?'', jawabannya ``Tidak apa-apa''
\end{itemize}

\textbf{Strengths:}
\begin{itemize}
    \item Sangat bisa diandalkan dalam krisis
    \item Tidak membuat drama yang tidak perlu
    \item Menghargai waktu sendiri pasangan
    \item Produktif dan fokus pada pekerjaan
\end{itemize}
\end{tipbox}

\section{Wave (Gelombang): Aku Tidak Bisa Tanpamu}

\subsection{Siapa itu Wave?}

Wave adalah tipe yang paling \textbf{emotional} dan \textbf{relational}. Mereka hidup untuk koneksi, tapi seringkali takut akan kehilangannya.

\begin{center}
\begin{tikzpicture}[
    wave/.style={draw=accent, line width=2pt, rounded corners=15pt, 
                 fill=accent!10, minimum width=10cm, inner sep=15pt}
]
    \node[wave] (box) {
        \begin{minipage}{9cm}
            \centering
            {\Large\bfseries\color{accent} KARAKTERISTIK WAVE}\\[0.5cm]
            \begin{itemize}[leftmargin=*, itemsep=0.2cm]
                \item \textbf{Generous \& Giving} --- pemberi yang murah hati
                \item \textbf{Focused on meaning} --- fokus pada makna dan emosi
                \item \textbf{Emotionally sensitive} --- sangat peka secara emosional
                \item \textbf{People-oriented} --- paling bahagia bersama orang lain
                \item \textbf{Clinging} --- mendekat saat stres (clinging group)
                \item \textbf{Ambivalent} --- kadang mendekat, kadang mendorong
                \item Takut \textbf{abandonment} (ditinggal) secara intens
                \item Sering merasa ``kurang'' atau tidak cukup baik
            \end{itemize}
        \end{minipage}
    };
\end{tikzpicture}
\end{center}

\subsection{Wave sebagai Anak}

Wave tumbuh dalam keluarga yang:
\begin{enumerate}
    \item Caregiver \textbf{membutuhkan perhatian anak} untuk merasa lebih baik (role reversal)
    \item Caregiver \textbf{kadang sangat perhatian, kadang tidak tersedia} (inconsistent)
    \item Saat anak sedih, caregiver \textbf{memfokuskan pada emosi sendiri}
    \item Caregiver \textbf{mudah tersinggung, sering kewalahan, atau menghukum}
    \item Caregiver \textbf{membuat anak sebagai ally} melawan pasangannya
    \item Saat anak ketakutan di malam hari, caregiver \textbf{tidak menenangkan dengan efektif}
\end{enumerate}

\subsection{Paradoks Wave: Mendekat tapi Mendorong}

Wave punya pola yang membingungkan:
\begin{enumerate}
    \item \textbf{Cling} --- mendekat, ingin perhatian, manja
    \item \textbf{Push away} --- tiba-tiba marah, menyindir, atau mengancam pergi
    \item \textbf{Test} --- ``Kalau aku pergi, apakah dia mengejar?''
    \item \textbf{Want pursuit} --- sebenarnya ingin pasangan mengejar dan membuktikan cinta
\end{enumerate}

\begin{praktikbox}[Mengapa Wave Melakukan Ini?]
Ini adalah \textbf{test of love}. Wave belajar dari pengalaman masa kecil bahwa cinta itu:
\begin{itemize}
    \item Tidak konsisten
    \item Harus diperjuangkan
    \item Bisa hilang kapan saja
\end{itemize}

Jadi mereka menciptakan situasi untuk membuktikan bahwa ``Dia benar-benar peduli'' dengan membuat pasangan \textbf{mengejar} mereka.
\end{praktikbox}

\section{Quiz Komprehensif: Kamu Tipe Apa?}

\begin{exercisebox}[Attachment Style Quiz --- Versi Detail]
\textbf{Petunjuk:} Untuk setiap pertanyaan, pilih jawaban yang \textit{paling sering} menggambarkan dirimu. Jujur pada dirimu sendiri --- ini untuk pemahaman, bukan penilaian.

\subsection*{BAGIAN A: Masa Kecil (Sebelum Usia 13)}

\textbf{1. Kapan kamu sedih atau menangis, bagaimana orang tua/caregiver bereaksi?}
\begin{itemize}
    \item[A] Datang, memeluk, menenangkan dengan efektif tanpa menghakimi
    \item[B] Memberikan mainan/uang, atau mengatakan ``Jangan menangis'', ``Besar kok nangis''
    \item[C] Kadang memeluk dengan penuh kasih, kadang bilang ``Ah, gitu aja nangis''
\end{itemize}

\textbf{2. Bagaimana budaya keluargamu tentang kebutuhan dan ketergantungan?}
\begin{itemize}
    \item[A] Normal dan diterima untuk saling bergantung dan dimanja
    \item[B] Ketergantungan dianggap kelemahan. ``Kamu harus mandiri''
    \item[C] Kamu diharapkan menjadi ``penyelamat'' atau tempat curhat orang tua
\end{itemize}

\textbf{3. Seberapa sering kamu merasa benar-benar ``dilihat'' dan dipahami oleh orang tua?}
\begin{itemize}
    \item[A] Sering --- mereka benar-benar tahu siapa aku
    \item[B] Jarang --- lebih fokus pada nilai/penampilanku
    \item[C] Kadang-kadang saja, tergantung mood mereka
\end{itemize}

\textbf{4. Saat kamu takut di malam hari dan memanggil, apa yang terjadi?}
\begin{itemize}
    \item[A] Seseorang datang dan menenangkan
    \item[B] Disuruh tidur sendiri, atau diberi lampu/mainan
    \item[C] Kadang ditemani, kadang disuruh ``jangan manja''
\end{itemize}

\textbf{5. Bagaimana cara orang tuamu menunjukkan kasih sayang?}
\begin{itemize}
    \item[A] Pelukan, ciuman, waktu berkualitas bersama
    \item[B] Hadiah, uang, atau prestasi yang dihargai
    \item[C] Kadang sangat manja, kadang dingin/marah
\end{itemize}

\subsection*{BAGIAN B: Dalam Hubungan Dewasa}

\textbf{6. Saat pasangan ingin bicara serius tentang hubungan, biasanya kamu:}
\begin{itemize}
    \item[A] Dengan senang hati --- ini memperkuat koneksi kita
    \item[B] Merasa ingin menjauh, menunda, atau menghilang
    \item[C] Senang, tapi takut dia akan marah besar atau pergi
\end{itemize}

\textbf{7. Bagaimana cara kamu ``recharge'' energi?}
\begin{itemize}
    \item[A] Bisa sama baiknya sendiri atau bersama pasangan
    \item[B] \textbf{Harus} punya waktu sendiri --- ``me-time'' adalah wajib
    \item[C] Butuh bersama pasangan --- sendiri membuatku cemas/ kosong
\end{itemize}

\textbf{8. Saat ada konflik, kamu cenderung:}
\begin{itemize}
    \item[A] Mencari solusi bersama, tetap terhubung meski marah
    \item[B] Menarik diri, merenung sendiri, atau menghindari
    \item[C] Mengejar, membuat dramatis, atau kadang mengancam putus
\end{itemize}

\textbf{9. Saat pasangan tidak membalas pesan dengan cepat, kamu:}
\begin{itemize}
    \item[A] Menganggap dia sibuk, tidak memikirkannya terus
    \item[B] Senang --- aku juga butuh ruang
    \item[C] Cemas, kepikiran terus, mungkin dia marah atau tidak cinta lagi
\end{itemize}

\textbf{10. Bagaimana perasaanmu tentang komitmen jangka panjang?}
\begin{itemize}
    \item[A] Nyaman dan positif --- aku ingin membangun sesuatu yang kuat
    \item[B] Was-was --- takut kehilangan kebebasan dan independensiku
    \item[C] Ingin sekali, tapi takut ditinggal atau tidak dicintai sepenuhnya
\end{itemize}

\subsection*{SCORING}

\begin{center}
\renewcommand{\arraystretch}{1.3}
\begin{tabular}{|c|c|c|c|}
\hline
\rowcolor{primary!20}
\textbf{Jawaban} & \textbf{Anchor} & \textbf{Island} & \textbf{Wave} \\
\hline
A & \checkmark & & \\
\hline
B & & \checkmark & \\
\hline
C & & & \checkmark \\
\hline
\end{tabular}
\end{center}

\textbf{Interpretasi:}
\begin{itemize}
    \item \textbf{7-10 A:} Kamu adalah \textbf{ANCHOR} yang kuat
    \item \textbf{7-10 B:} Kamu adalah \textbf{ISLAND} yang kuat
    \item \textbf{7-10 C:} Kamu adalah \textbf{WAVE} yang kuat
    \item \textbf{Campuran:} Sangat normal! Sebut saja \textbf{Anchor-ish}, \textbf{Island-ish}, atau \textbf{Wave-ish}
\end{itemize}

\textit{Catatan: Tidak ada yang salah dengan tipe apapun. Yang penting adalah kesadaran dan kesediaan untuk berkembang.}
\end{exercisebox}

\section{Sherlocking: Menyelidiki Gaya Pasanganmu}

Setelah memahami dirimu sendiri, langkah krusial berikutnya adalah memahami pasanganmu. Dr. Tatkin menyebut ini \keyterm{Sherlocking}.

\subsection{Pilar Sherlock yang Efektif}

\begin{tipbox}[5 Pilar Sherlocking]
\begin{enumerate}
    \item \textbf{Observasi Non-Verbal}
    \begin{itemize}
        \item Mata: kontak mata, pupil, arah pandangan
        \item Wajah: ekspresi, simetri, warna
        \item Tubuh: postur, gestur, ketegangan otot
        \item Suara: nada, kecepatan, volume
    \end{itemize}
    
    \item \textbf{Dengar Antara Baris}
    \begin{itemize}
        \item Apa yang \textit{ditonjolkan} vs apa yang \textit{disederhanakan}?
        \item Apa yang \textit{dihindari}?
        \item Pola cerita yang berulang?
    \end{itemize}
    
    \item \textbf{Tanyakan dengan Lembut}
    \begin{itemize}
        \item ``Ceritakan tentang keluargamu waktu kecil''
        \item Bukan interogasi, tapi rasa ingin tahu
        \item Bagikan pengalamanmu dulu untuk membuka percakapan
    \end{itemize}
    
    \item \textbf{Amati dalam Berbagai Situasi}
    \begin{itemize}
        \item Saat stres: apa respons pertamanya?
        \item Saat senang: bagaimana cara berbagi?
        \item Saat marah: apa coping mechanism-nya?
        \item Dengan orang lain vs denganmu
    \end{itemize}
    
    \item \textbf{Catat Pola}
    \begin{itemize}
        \item Apa yang memicu penarikan diri?
        \item Apa yang membuatnya merasa aman?
        \item Bagaimana dia meminta koneksi (meski tidak langsung)?
    \end{itemize}
\end{enumerate}
\end{tipbox}

\subsection{Tabel Referensi Cepat}

\begin{center}
\renewcommand{\arraystretch}{1.6}
\begin{tabular}{|>{\raggedright}p{2.8cm}|>{\raggedright}p{4.5cm}|>{\raggedright}p{4.5cm}|>{\raggedright\arraybackslash}p{4.5cm}|}
\hline
\rowcolor{primary!20}
\textbf{Aspek} & \textbf{Anchor} & \textbf{Island} & \textbf{Wave} \\
\hline
\textbf{Cara bicara} & Koheren, ekspresif, jelas & Ringkas, kadang terputus, logis & Banyak, emosional, hiperbolis, ``selalu''/``tidak pernah'' \\
\hline
\textbf{Kontak mata} & Nyaman, hangat, stabil & Sering menghindar, lirik sebentar & Intens, menuntut, atau menghindar saat marah \\
\hline
\textbf{Sentuhan fisik} & Natural, tepat waktu, nyaman & Canggung, terukur, batasan jelas & Banyak, memerlukan, tapi kadang menolak \\
\hline
\textbf{Saat stres} & Mencari solusi bersama, tetap terbuka & Menjauh, merenung sendiri, ``butuh ruang'' & Mengejar, emosional, atau marah dulu \\
\hline
\textbf{Konflik} & Kolaboratif, repair cepat & Menghindar atau menyerang balik & Emosional, dramatis, kadang mengancam putus \\
\hline
\textbf{Kebutuhan utama} & Seimbang antara kedekatan \& ruang & Banyak ruang sendiri, jangan dipaksa & Banyak kedekatan, reassurance, jangan ditinggal \\
\hline
\textbf{Ketakutan terbesar} & Kehilangan koneksi yang sudah dibangun & Kehilangan otonomi, terjebak, tenggelam & Ditinggal, tidak dicintai, dikhianati \\
\hline
\textbf{Cara meminta kasih sayang} & Langsung, jelas, tanpa malu & Tidak meminta --- ``aku fine'', terkejut saat ditanya & Kadang tidak langsung (nge-test), kadang terlalu banyak \\
\hline
\textbf{Memori masa kecil} & Jelas, umumnya positif & Kabur, sering ``positif'' tapi minim detail & Campur aduk, intens, mungkin trauma \\
\hline
\end{tabular}
\end{center}

\section{Mengapa Memahami Ini Menyelamatkan Hubunganmu}

\subsection{Mengurangi Personifikasi}

Tanpa pemahaman ini:
\begin{itemize}
    \item ``Dia tidak membalas chat karena tidak cinta lagi padaku!''
    \item ``Dia selalu menghindar saat ada masalah, dia tidak dewasa!''
    \item ``Dia terlalu manja, dia lemah!''
\end{itemize}

Dengan pemahaman ini:
\begin{itemize}
    \item ``Dia tidak membalas chat --- mungkin dia Island yang butuh ruang untuk recharge. Ini bukan tentangku.''
    \item ``Dia menghindar --- itu coping mechanism dari masa kecilnya. Bagaimana aku bisa membuatnya merasa aman?''
    \item ``Dia butuh banyak reassurance --- itu karena takut ditinggal. Bagaimana aku bisa meyakinkannya?''
\end{itemize}

\begin{insightbox}
\textbf{Paradigma Shift:} Konflik bukan lagi ``dia vs aku'', tapi ``kita vs perbedaan wiring kita''.
\end{insightbox}

\subsection{Menyesuaikan Approach}

\begin{praktikbox}[Cara Menghadapi Masing-Masing Tipe]
\textbf{Jika Pasanganmu Anchor:}
\begin{itemize}
    \item Bagus! Dia paling mudah diajak kerja sama
    \item Tetap lakukan maintenance --- jangan anggap remeh
    \item Gunakan sebagai ``secure base'' untuk tumbuh bersama
\end{itemize}

\textbf{Jika Pasanganmu Island:}
\begin{itemize}
    \item Jangan \textbf{crowd} (terlalu mendesak) --- beri ruang
    \item Jangan ambil penarikan diri secara personal
    \item Beri reassurance: ``Aku akan selalu di sini''
    \item Ajak bicara saat dia \textit{sudah} recharge, bukan saat dia menarik diri
    \item Hindari kesan ``needy'' atau ``controlling''
\end{itemize}

\textbf{Jika Pasanganmu Wave:}
\begin{itemize}
    \item Jangan biarkan dia merasa ditinggal
    \item Beri banyak reassurance verbal dan fisik
    \item Saat dia marah/withdraw, \textbf{kejar dengan lembut} (gentle pursuit)
    \item Validasi perasaannya sebelum memberi solusi
    \item Jangan bilang ``tenang'' atau ``jangan terlalu sensitif''
\end{itemize}
\end{praktikbox}

\section{Moving Toward Secure Functioning}

\begin{insightbox}
\textbf{Key Takeaway:} Tidak peduli tipe aslimu apa, \textbf{setiap orang} bisa bergerak menuju \textit{secure functioning}. Ini bukan tentang mengubah dirimu jadi orang lain, tapi tentang:
\begin{enumerate}
    \item \textbf{Kesadaran} akan wiring-mu
    \item \textbf{Kesediaan} untuk bertanggung jawab atas reaksi-mu
    \item \textbf{Kesepakatan} dengan pasangan untuk saling mendukung
\end{enumerate}
\end{insightbox}

\subsection{Apa yang Bisa Dilakukan Setiap Tipe}

\begin{center}
\renewcommand{\arraystretch}{1.4}
\begin{tabular}{|>{\raggedright}p{3cm}|>{\raggedright\arraybackslash}p{12cm}|}
\hline
\rowcolor{primary!20}
\textbf{Tipe} & \textbf{Action Items} \\
\hline
\textbf{Anchor} & 
\begin{itemize}[leftmargin=*, itemsep=0.1cm]
    \item Terus jadi ``secure base'' untuk pasangan
    \item Jangan jenuh dengan maintenance hubungan
    \item Ajari pasanganmu tentang konsep ini
\end{itemize} \\
\hline
\textbf{Island} & 
\begin{itemize}[leftmargin=*, itemsep=0.1cm]
    \item Sadari bahwa ``menjauh'' bukan solusi jangka panjang
    \item Latih diri untuk meminta bantuan --- ini bukan kelemahan
    \item Beritahu pasangan: ``Aku butuh 30 menit sendiri, tapi aku akan kembali''
    \item Izinkan pasangan merawatmu
\end{itemize} \\
\hline
\textbf{Wave} & 
\begin{itemize}[leftmargin=*, itemsep=0.1cm]
    \item Sadari bahwa ``testing'' bisa melukai pasangan
    \item Latih self-soothing --- kamu bisa menenangkan diri sendiri juga
    \item Komunikasikan kebutuhan langsung tanpa drama
    \item Terima reassurance tanpa meragukan kejujurannya
\end{itemize} \\
\hline
\end{tabular}
\end{center}

\section{Action Steps untuk Bab Ini}

\begin{exercisebox}[Latihan Bersama Pasangan]
\textbf{Langkah 1: Individual Assessment (15 menit)}
\begin{itemize}
    \item Masing-masing kerjakan quiz secara terpisah
    \item Tuliskan hasil dan refleksi singkat
\end{itemize}

\textbf{Langkah 2: Share Session (30-45 menit)}
\begin{itemize}
    \item Duduk berhadapan, pegang tangan
    \item Bergantian sharing hasil quiz
    \item Jangan defensive --- ini bukan kritik, tapi pemahaman
    \item Contoh kalimat: ``Aku melihat diriku sebagai Island karena...''
\end{itemize}

\textbf{Langkah 3: Discussion Questions}
\begin{enumerate}
    \item ``Bagaimana tipe kita saling melengkapi?''
    \item ``Bagaimana kita bisa saling mendukung saat...''
    \begin{itemize}
        \item Aku butuh ruang (jika ada Island)
        \item Aku cemas (jika ada Wave)
    \end{itemize}
    \item ``Apa yang bisa kita lakukan berbeda mulai sekarang?''
\end{enumerate}

\textbf{Langkah 4: Buat Agreement}
\begin{itemize}
    \item Tulis 3 komitmen spesifik berdasarkan pemahaman ini
    \item Contoh: ``Jika aku (Island) butuh ruang, aku akan bilang 'Aku butuh 1 jam sendiri, aku akan balik pukul 8' daripada menghilang''
\end{itemize}
\end{exercisebox}
